% Local compensation for the blank line inserted after the abstract.
% Applies only to this opening section; global heading spacing is unchanged.
\vspace*{-30pt}
% Preserve the existing page break so subsequent pages remain unchanged.
\enlargethispage{-30pt}
\section{问题重述}

\subsection{问题背景}

药材烘干的目标并非简单地使表面水分降低，而是在不过度加热和不过度干燥表层的前提下，使药材内部所有位置均达到规定含水标准。烘房温度升高后，热量由外向内传递，表层水分同时在内部浓度梯度和表面对流传质的共同作用下向外迁移。由于传热和传质的特征时间不同，药材内部会长期存在“表面已干、中心仍湿”的空间不均匀现象。因此，只记录环境状态或表面状态不能直接判断整根药材是否完成烘干，必须建立能够描述内部时空分布的非稳态模型。

本题以近似圆柱形药材为研究对象，其长为 25~cm、初始半径为 2~cm，初始温度和水分浓度分别为 $28~{}^\circ\mathrm C$ 和 $2.55~\mathrm{kg/kg}$。烘干过程中，外部温湿环境随时间变化，材料热物性和水分扩散能力也会随状态改变，并伴随明显的尺寸收缩。由此，模型不仅需要刻画固定几何条件下的热湿传递，还需进一步处理变物性、全截面达标判定及移动边界等问题。四问正是沿着这一物理复杂性逐步展开，形成由基础传热传质模型向非线性收缩模型递进的建模过程。

\subsection{问题提出}

根据题目给定的几何尺寸、初始状态、环境条件及物性关系，需要建立药材烘干过程中的热湿迁移模型，获得温度和水分浓度的时空分布，并进一步确定药材达到规定含水标准的烘干时间。随着物性变化和尺寸收缩的引入，模型还需逐步处理非线性耦合与移动边界问题。

\textbf{问题一：}首先需要刻画预热阶段药材内部热量与水分的径向迁移规律。由于药材内部存在明显的温度和水分梯度，应建立圆柱坐标下的非稳态传热传质模型，并结合时变环境边界求得温湿场，为后续各问提供统一的基础模型和数值求解框架。

\textbf{问题二：}在基础模型上进一步考虑物性随温度和水分状态变化的影响，使导热、储热和水分扩散过程相互耦合。由此需将问题一的基准模型扩展为变物性的非线性热湿迁移模型，以分析整个烘干过程中温度场与水分场的动态演化及其控制机制。

\textbf{问题三：}在获得长期水分迁移规律后，需要进一步解决“何时完成烘干”的判定问题。将“药材各处水分均低于规定阈值”转化为空间最大水分浓度的首达条件，并据此确定控制位置及最终烘干时间，从而避免仅依据表面或平均水分作出过早停机判断。

\textbf{问题四：}实际烘干过程中药材半径持续收缩，使计算区域由固定域转化为移动域。为此需在变物性热湿模型中进一步引入时变几何边界，通过固定域映射处理收缩过程，并分析扩散距离缩短对烘干效率的影响；同时采用独立移动网格计算以及COMSOL物理仿真对移动边界结果进行交叉验证。
